\begin{abstract}
Federated LoRA comparisons often rest on a single non-IID partition draw, even when reported method gaps sit near ordinary run-to-run noise. We ask how much of the variation in held-out next-token loss comes from the partition draw, how often few-draw rankings flip, and how many draws a target effect size requires. In a pre-registered design of $120$ runs at two scales (TinyLlama-1.1B-Chat and LLaMA-3.2-3B), with partition seeds separated from training seeds, we fit a two-way variance decomposition for FedIT, FFA-LoRA, and FLoRA on Dolly, partitioned by its category labels at Dirichlet $\alpha{=}0.1$. On TinyLlama the partition and interaction shares are $0.508$ and $0.469$ against a residual of $0.023$. On LLaMA-3.2-3B the interaction dominates ($0.764$) and the partition interval includes zero. Ranking flips concentrate in the near-tied FedIT--FLoRA pair: the three-draw paired flip probability is $0.377$, while gaps of about $0.023$ never reverse in a single draw. Detecting that observed near-tie at $80\%$ power needs $206$ paired partition draws on TinyLlama and $46$ on LLaMA-3.2-3B; separable pairs need only $3$. We recommend reporting the number of partition draws, pairing comparisons on partitions, and stating the smallest difference the design can detect. We make no method recommendation.
\end{abstract}

\begin{IEEEkeywords}
federated learning; low-rank adaptation; non-IID partitioning; experimental variance; reproducibility; large language models
\end{IEEEkeywords}
